\chapter*{Abstract}
\setheader{Abstract}

Diese Bachelorarbeit beschäftigt sich mit der Weiterentwicklung der Check-API und das damit verbundene System der Programmierlernplattform CodeTask. CodeTask ist eine an der HTWG Konstanz entwickelte Website zum bearbeiten, lösen und lernen von Programmieraufgaben in der Programmiersprache Scala. CodeTask ermöglicht dabei direkte Ausführung und Fehlerbehandlung von Code. Für diese Aufgabe kommt der Isolierte Service Check-API in einsatz, welche darauf ausgelegt ist Code in einem aktiven System ausführen zu können. Zu Beginn der Arbeit verfügte die Check-API bereit über die Fähigkeit Code zusammen mit öffenltichen Tests ausführen und evaluieren zu können. Die Check-API verfügte zu diesem Zeitpunkt allerdings nur über eine eingeschränkte Rückmeldung, besaß keine Unterstützung für versteckte Tests, war nur eingeschränkt gegen problematischen oder bösartige Einreichungen abgesichert und konnte durch eine einfache unendliche Schleife zum Stillstand kommen

Ziel der Arbeit war es, die Check-API robuster, aussagekräftiger und sicherer zu machen. Dafür wurde das System um die Möglichkeit erweitert versteckte Tests zu nutzen, sodass gezielte Umgehungen der öffentlihen Tests erschwert werden. Zusätzlich wurde die Check-API mit einer sicherheitsprüfung von Code auf unsichere Befehle ausgestattet, um bekannte kritische Befehle vor der Ausführung zu erkennen und ausgewählte Angrffs- bzw. Fehlerfälle zu verhindern. Zudem wurden Fehlermeldungen systematisch ausgewertet und aufbereitet um eine verbesserte Fehlersuche für den Nutzer zu ermöglichen.

Die Evaluation zeigt, dass in den getesteten Fällen keine inhalte der versteckten Tests an Nutzer zurückgegeben wurden. Des weiteren werden Compilierungsfehler, fehlgeschlagene öffentliche, fehlgeschlagene versteckte Tests, Zeitüberschreitungen und Sicherheitsverstöße korrekt erkannt, verarbeitet und strukturiert an den Nutzer weitergegeben. Die Performance-Messungen legen eine klare Verbesserung der Leistung offen, zeigen jedoch deutlich, dass die Ausführung von Nutzercode rechenintensiv bleibt und für hohe Last geeignete Skalierungsmaßnahmen erfoderlich sind.